% !TeX root = ../main.tex
\section{Parsing}

The lexer hands you a flat stream of tokens. That stream has no structure --- it's just a
list. \texttt{IDENT(x) ASSIGN(=) NUMBER(3) PLUS(+) NUMBER(4)} doesn't say, by itself,
whether this is one statement or two, or which operator applies to which operand. The
\textbf{parser}'s job is to take that stream and, using the grammar you wrote in the
previous block, recover the structure the programmer intended: a tree. The next section
formalizes exactly what that tree looks like for a small imperative language (assignments,
sequencing, conditionals, loops); here we care about the mechanics of \emph{building} it,
whatever shape it ends up taking.

If a token sequence can't be produced by the grammar at all, the parser reports a
\textbf{syntax error} --- a different failure mode from the lexer's ``I don't recognize
this character'' errors. \texttt{x = = 3} tokenizes just fine (every character is legal),
but \texttt{ASSIGN ASSIGN NUMBER} isn't a valid statement in any reasonable grammar.

\subsection{recursive descent: one function per nonterminal}

The most direct way to write a parser by hand is \textbf{recursive descent}: for every
nonterminal in the grammar, write a function; that function's job is to recognize (and
consume) exactly the tokens that nonterminal can produce, calling other nonterminals'
functions as it goes. It's a literal, mechanical translation from grammar to code.

Take a small arithmetic grammar:
\[
\begin{aligned}
E  &::= T\; E' \\
E' &::= \texttt{+}\; T\; E' \mid \texttt{-}\; T\; E' \mid \varepsilon \\
T  &::= F\; T' \\
T' &::= \texttt{*}\; F\; T' \mid \texttt{/}\; F\; T' \mid \varepsilon \\
F  &::= \texttt{NUMBER} \mid \texttt{(}\; E\; \texttt{)}
\end{aligned}
\]
This grammar bakes precedence directly into its structure: $T$ (terms, built from
\texttt{*}/\texttt{/}) sits ``inside'' $E$ (expressions, built from \texttt{+}/\texttt{-}),
so a $T$ always binds tighter. Translating it to recursive descent is almost mechanical:

\begin{lstlisting}[language=Python]
def E():
    T()
    Eprime()

def Eprime():
    if peek() in ("+", "-"):
        advance()
        T()
        Eprime()
    # else: epsilon. FOLLOW(E') = {')', EOF} tells us it's safe to
    # stop here without consuming anything.

def T():
    F()
    Tprime()

def Tprime():
    if peek() in ("*", "/"):
        advance()
        F()
        Tprime()

def F():
    if peek() == "NUMBER":
        advance()
    elif peek() == "(":
        advance()
        E()
        expect(")")
    else:
        error("expected NUMBER or '('")
\end{lstlisting}

Every production became a function; every alternative in the grammar became a branch
picked by looking at the next token.

\subsection{FIRST/FOLLOW sets and what LL(1) means}

That last function, \texttt{F}, decides which branch to take by looking at exactly one
token of input. That's not an accident --- it's a property of the grammar, and it's called
\textbf{LL(1)}: parseable Left-to-right, building a Leftmost derivation, with 1 token of
lookahead.

\begin{definition}[FIRST set]
$\mathrm{FIRST}(\alpha)$ is the set of terminals that can appear as the first symbol of
some string derived from $\alpha$ (plus $\varepsilon$, if $\alpha$ can derive the empty
string).
\end{definition}

\begin{definition}[FOLLOW set]
$\mathrm{FOLLOW}(A)$ is the set of terminals that can appear immediately after $A$ in some
derivation from the start symbol (plus the end-of-input marker, if $A$ can be the last
symbol of a full derivation).
\end{definition}

For the grammar above:
\[
\mathrm{FIRST}(E) = \mathrm{FIRST}(T) = \mathrm{FIRST}(F) = \{\texttt{NUMBER}, \texttt{(}\}
\qquad
\mathrm{FIRST}(E') = \{\texttt{+}, \texttt{-}, \varepsilon\}
\]
\[
\mathrm{FOLLOW}(E) = \{\texttt{)}, \mathtt{EOF}\}
\qquad
\mathrm{FOLLOW}(E') = \mathrm{FOLLOW}(E)
\]

A grammar is \textbf{LL(1)} when, for every nonterminal with more than one alternative,
looking at just the next token (comparing it against the FIRST sets of each alternative,
and FOLLOW when $\varepsilon$ is in play) always tells you unambiguously which alternative
to take. That's precisely what lets \texttt{Eprime} decide, with a single \texttt{peek()},
whether to recurse or to stop: if the next token is in $\{\texttt{+}, \texttt{-}\}$ take
the first branch, and if it's in $\mathrm{FOLLOW}(E') = \{\texttt{)}, \mathtt{EOF}\}$ take
the $\varepsilon$ branch. If those sets ever overlapped, one token of lookahead wouldn't be
enough, and plain recursive descent (without backtracking) couldn't handle the grammar
as written.

\subsection{grammar massaging: left recursion and left factoring}

Not every grammar you'd naturally write is LL(1), and recursive descent doesn't work at
all on some of them --- it just hangs. Two transformations fix the common cases.

\textbf{Left recursion.} The grammar $E ::= E\; \texttt{+}\; T \mid T$ is the natural way
to write ``addition is left-associative,'' but translate it directly into a function and
you get:
\begin{lstlisting}[language=Python]
def E():
    E()      # infinite recursion before consuming a single token
    expect("+")
    T()
\end{lstlisting}
\texttt{E} calls itself before consuming any input, forever. The fix is exactly the
transformation that produced the grammar we used above: rewrite
$E ::= E\;\texttt{+}\;T \mid T$ as $E ::= T\;E'$, $E' ::= \texttt{+}\;T\;E' \mid
\varepsilon$. This is a completely general rewrite (any left-recursive rule
$A ::= A\alpha \mid \beta$ becomes $A ::= \beta A'$, $A' ::= \alpha A' \mid \varepsilon$)
and it preserves the language while turning a self-call at the start of a production into
a self-call after consuming a token --- which is exactly what a loop instead of unbounded
recursion looks like.

\textbf{Left factoring.} Two alternatives that start with the same token also break
LL(1), even without recursion. Consider statements:
\[
\mathit{Stmt} ::= \texttt{if}\;(\mathit{E})\;\mathit{Stmt}
\ \mid\ \texttt{if}\;(\mathit{E})\;\mathit{Stmt}\;\texttt{else}\;\mathit{Stmt}
\]
Both alternatives start with \texttt{if (E) Stmt}; peeking at the next token after that
tells you nothing until you've already parsed the whole \texttt{if}-body. Left factoring
pulls the common prefix into its own rule and defers the decision:
\[
\mathit{Stmt} ::= \texttt{if}\;(\mathit{E})\;\mathit{Stmt}\;\mathit{ElseOpt}
\qquad
\mathit{ElseOpt} ::= \texttt{else}\;\mathit{Stmt} \mid \varepsilon
\]
Now the choice (\texttt{else} or not) is made by a single, separate function that peeks at
one token, exactly like \texttt{Eprime} above --- and as a bonus this is precisely the
grammar shape that resolves the classic ``dangling else'' ambiguity in favor of attaching
\texttt{else} to the nearest unmatched \texttt{if}.

\subsection{Pratt parsing: precedence climbing for expressions}

Recursive descent works, but writing a nonterminal per precedence level ($E$, $T$, $F$,
and one more layer for every new operator) gets tedious fast, especially once you add
exponentiation, comparisons, boolean \texttt{and}/\texttt{or}, and so on. \textbf{Pratt
parsing} (also called \textbf{precedence climbing} or \textbf{operator-precedence
parsing}) replaces that tower of nonterminals with a single loop driven by a table of
numeric \emph{binding powers}.

The idea: parse one ``atom'' (a number or a parenthesized expression), then keep looking
at the next operator. If its binding power is at least as strong as the minimum we're
currently allowed to accept, consume it, recursively parse the right-hand side (demanding
one more unit of binding power, to force left-associativity), and fold the result into the
left-hand side. Stop the moment the next operator is too weak, or there isn't one.

\begin{lstlisting}[language=Python]
PRECEDENCE = {"+": 1, "-": 1, "*": 2, "/": 2}

def parse_expr(tokens, min_bp=0):
    left = parse_atom(tokens)
    while True:
        op = tokens.peek()
        if op not in PRECEDENCE or PRECEDENCE[op] < min_bp:
            break
        tokens.advance()
        bp = PRECEDENCE[op]
        right = parse_expr(tokens, bp + 1)   # +1: left-associative
        left = ("BinOp", op, left, right)
    return left

def parse_atom(tokens):
    tok = tokens.advance()
    if tok == "(":
        e = parse_expr(tokens, 0)
        tokens.expect(")")
        return e
    return ("Num", tok)
\end{lstlisting}

Adding a new precedence level is now a one-line change to a dictionary, not a new
function and a new grammar rule. This is the standard trick used in real
implementations (it's how Pratt originally described it in 1973, and it's what most
hand-written expression parsers use today, including the one built step by step in
Nystrom's \emph{Crafting Interpreters}).

\begin{example}[Tracing precedence climbing on \texttt{1 + 2 * 3}]
Call \texttt{parse\_expr(tokens, min\_bp=0)}:
\begin{enumerate}
  \item \texttt{left = parse\_atom()} $\to$ \texttt{Num(1)}.
  \item Peek \texttt{+}: precedence 1 $\ge$ 0, so consume it. Recurse with
    \texttt{min\_bp = 1 + 1 = 2}:
    \begin{enumerate}
      \item \texttt{left = parse\_atom()} $\to$ \texttt{Num(2)}.
      \item Peek \texttt{*}: precedence 2 $\ge$ 2, so consume it. Recurse with
        \texttt{min\_bp = 2 + 1 = 3}:
        \begin{itemize}
          \item \texttt{left = parse\_atom()} $\to$ \texttt{Num(3)}. No more tokens, return
            \texttt{Num(3)}.
        \end{itemize}
      \item \texttt{left} becomes \texttt{BinOp(*, Num(2), Num(3))}. No more tokens,
        return.
    \end{enumerate}
  \item \texttt{right} above is \texttt{BinOp(*, Num(2), Num(3))}. \texttt{left} becomes
    \texttt{BinOp(+, Num(1), BinOp(*, Num(2), Num(3)))}.
\end{enumerate}
The result is $1 + (2 \times 3)$, not $(1+2)\times 3$ --- multiplication's higher binding
power forced it to bind to \texttt{2} before control returned to the outer \texttt{+}, with
no explicit grammar rule for precedence anywhere in sight.
\end{example}

\subsection{bottom-up parsing, briefly}

Recursive descent and Pratt parsing are both \textbf{top-down}: they start from the start
symbol and work their way down to the tokens, guessing which production to expand based on
lookahead. The other family, \textbf{LR (shift-reduce) parsing}, works \textbf{bottom-up}:
it pushes tokens onto a stack (``shift'') and, whenever the top of the stack matches the
right-hand side of some production, replaces it with the corresponding nonterminal
(``reduce''), building the tree from the leaves up to the root.

LR parsers accept a strictly larger class of grammars than LL --- in particular, they
handle left recursion natively and don't require left factoring. The price is that the
parsing tables (which state does what on which lookahead token) are large and genuinely
unpleasant to construct by hand; in practice nobody does. Tools like \texttt{yacc},
\texttt{bison}, and ANTLR take a grammar and generate the table (and often the whole
parser) for you. Table construction itself (SLR, LALR, canonical LR) is out of scope for
these notes --- Aho, Lam, Sethi and Ullman's \emph{Compilers: Principles, Techniques, and
Tools} devotes several chapters to it if you want the full theory. For hand-written
parsers, which is the practical choice for most small languages and DSLs, top-down wins
on readability and error-message quality, and that's what the rest of this block assumes.

\subsection{error recovery, briefly}

A parser that stops dead at the first syntax error is unpleasant to use --- you fix one
typo, rerun, and find the \emph{next} one, one at a time. The standard fix by hand is
\textbf{panic mode}: on an error, discard tokens until you reach a known
\emph{synchronization point} (a semicolon, a keyword that starts a new statement, a
closing brace) and resume parsing from there, so the rest of the file still gets checked.
This is exactly the strategy discussed in more depth, with a concrete DSL example, in the
``Error Recovery'' part of \texttt{topics/06-dsl-design/sections/02-design-decisions.tex}
--- worth a look if you're implementing a real parser and want the fuller design
discussion; here the point is just to know the technique exists and what it's called.

\subsection{a worked example: parsing \texttt{3 + 4 * 2 - 1} by hand}

Using the LL(1) grammar from the top of this section (not precedence climbing, so you can
see the recursive-descent call structure explicitly), trace the token stream
\texttt{NUMBER(3) PLUS NUMBER(4) STAR NUMBER(2) MINUS NUMBER(1) EOF}.

\begin{enumerate}
  \item \texttt{E()} calls \texttt{T()}, which calls \texttt{F()}: consumes \texttt{3}.
    Back in \texttt{T()}, call \texttt{Tprime()}: peek is \texttt{+}, not in
    $\{\texttt{*},\texttt{/}\}$, so \texttt{Tprime} takes the $\varepsilon$ branch. \texttt{T}
    returns having matched just \texttt{3}.
  \item Back in \texttt{E()}, call \texttt{Eprime()}: peek is \texttt{+}, consume it, call
    \texttt{T()} again.
  \item This inner \texttt{T()} calls \texttt{F()}: consumes \texttt{4}. Then
    \texttt{Tprime()}: peek is \texttt{*}, consume it, call \texttt{F()}: consumes
    \texttt{2}. \texttt{Tprime()} recurses: peek is \texttt{-}, not in
    $\{\texttt{*},\texttt{/}\}$, take $\varepsilon$. This \texttt{T} has matched
    \texttt{4 * 2}.
  \item Back in the outer \texttt{Eprime()}, having matched \texttt{+ 4 * 2}, recurse into
    \texttt{Eprime()} again: peek is \texttt{-}, consume it, call \texttt{T()}, which
    matches just \texttt{1} (its own \texttt{Tprime} immediately sees \texttt{EOF} and
    takes $\varepsilon$).
  \item \texttt{Eprime()} recurses one more time: peek is \texttt{EOF}, take
    $\varepsilon$, unwind everything.
\end{enumerate}

Reading the calls in the order they combined values (not the order they were entered)
gives the tree, using the same node-drawing convention the next section formalizes for
full programs:
\[
\begin{array}{c}
\textsc{Minus} \\
\swarrow \quad \searrow \\
\textsc{Plus} \qquad 1 \\
\swarrow \quad \searrow \\
3 \qquad \textsc{Star} \\
\quad\quad\quad \swarrow \ \searrow \\
\quad\quad\quad 4 \quad\ \ 2
\end{array}
\]
which is exactly $(3 + (4 \times 2)) - 1$ --- the grammar's built-in precedence
(multiplication nested one level deeper than addition/subtraction) produced the right
tree without any explicit binding-power bookkeeping, at the cost of four nonterminals
instead of one loop.

\subsection{Exercises}
\begin{itemize}
  \item Eliminate left recursion from
    $\mathit{Cmp} ::= \mathit{Cmp}\ \texttt{<}\ E \mid \mathit{Cmp}\ \texttt{>}\ E \mid E$,
    then write the corresponding recursive-descent functions.
  \item Compute $\mathrm{FIRST}$ and $\mathrm{FOLLOW}$ for every nonterminal in
    $S ::= \texttt{if}\ (E)\ S\ \mathit{ElseOpt} \mid \texttt{print}\ E$,
    $\mathit{ElseOpt} ::= \texttt{else}\ S \mid \varepsilon$. Is this grammar LL(1)?
  \item Trace the precedence-climbing code on \texttt{2 * 3 + 4 * 5}, showing the sequence
    of recursive calls and the final tree, the way the worked example does for
    \texttt{1 + 2 * 3}.
  \item Extend \texttt{PRECEDENCE} and \texttt{parse\_expr} to support a right-associative
    exponentiation operator \texttt{\^{}} (binding tighter than \texttt{*}/\texttt{/}).
    Hint: right-associativity means the recursive call on the right should \emph{not} add
    1 to the minimum binding power.
  \item Using the LL(1) grammar from this section, hand-parse
    \texttt{(1 + 2) * 3} step by step and draw the resulting tree.
  \item In your own words: why does left recursion break plain recursive descent but pose
    no problem at all for a bottom-up (shift-reduce) parser?
\end{itemize}
